\chapter{1979年全国卷（理）}
\section{解答题}

\begin{problem}
若 \(\left(z-x\right)^2-4\left(x-y\right)\left(y-z\right)=0\)，求证：\(x\)，\(y\)，\(z\) 成等差数列.
\end{problem}

\begin{solution}
令
\(
a=x-y,\quad b=y-z.
\)
则
\(
z-x=-(a+b).
\)
代入已知等式，得到
\[
(a+b)^2-4ab=0
\]

展开左边并因式分解，得
\[
 (a-b)^2=0
\]

实数的平方等于零只能说明该实数本身等于零，所以
\(
a=b
\)，即
\[
x-y=y-z
\]
这正是数列 \(x\)，\(y\)，\(z\) 成等差数列的判定条件，故命题得证.
\end{solution}

\begin{problem}
化简： \(\frac{1}{1-\frac{1}{1-\frac{1}{1-\csc^2 x}}}\).
\end{problem}

\begin{solution}
设
\(
u=\csc^2x.
\)
首先确定原式的定义域. 必须有 \(\sin x\ne0\)，并且最内层分式的分母
\(1-\csc^2x\ne0\)，即 \(u\ne1\). 因为 \(u=1/\sin^2x\geq1\)，排除 \(u=1\) 后有 \(u>1\). 这两个条件等价于
\(
x\notin\frac{\pi}{2}\Z.
\)

在这个定义域内，先化简最内层的下一层分母：
\[
1-\frac{1}{1-u}=\frac{1-u-1}{1-u}=\frac{-u}{1-u}=\frac{u}{u-1}
\]

由于 \(u>1\)，这个分母不为零. 因此原式最外层分母为
\[
1-\frac{1}{1-\frac{1}{1-u}}
=1-\frac{1}{u/(u-1)}
=1-\frac{u-1}{u}
=\frac{1}{u}
\]

于是原式的值为
\[
\frac{1}{1/u}=u=\csc^2x
\]
所以化简结果为 \(\csc^2x\)，但必须附带原式的定义域
\(
x\notin\frac{\pi}{2}\Z
\)
（\(\Z\) 表示整数集）.
\end{solution}

\begin{problem}
甲，乙二容器内都盛有酒精，甲有 \(V_{1}\text{ 公斤}\)，乙有 \(V_{2}\text{ 公斤}\). 甲中纯酒精与水（重量）之比为 \(m_{1}:n_{1}\)，乙中纯酒精与水之比为 \(m_{2}:n_{2}\). 问将二者混合后所得液体中纯酒精与水之比是多少？
\end{problem}

\begin{solution}
把每个容器中的总重量看作相应比例的总份数. 甲中纯酒精与水的重量比为 \(m_1:n_1\)，所以甲中纯酒精和水的重量分别为
\(\frac{m_1V_1}{m_1+n_1}\)，\(\frac{n_1V_1}{m_1+n_1}\).
同理，乙中纯酒精和水的重量分别为
\(\frac{m_2V_2}{m_2+n_2}\)，\(\frac{n_2V_2}{m_2+n_2}\).

混合时两种物质的重量分别相加. 因此混合后纯酒精与水的重量之比为
\[
\left(\frac{m_{1}V_{1}}{m_{1}+n_{1}}+\frac{m_{2}V_{2}}{m_{2}+n_{2}}\right):\left(\frac{n_{1}V_{1}}{m_{1}+n_{1}}+\frac{n_{2}V_{2}}{m_{2}+n_{2}}\right)
\]

比例两项可以同时乘以正数 \(\left(m_1+n_1\right)\left(m_2+n_2\right)\)，不改变比例值，故所求之比为
\[
\bigl[m_{1}V_{1}(m_{2}+n_{2})+m_{2}V_{2}(m_{1}+n_{1})\bigr]:\bigl[n_{1}V_{1}(m_{2}+n_{2})+n_{2}V_{2}(m_{1}+n_{1})\bigr]
\]
\end{solution}

\begin{problem}
叙述并证明勾股定理.
\end{problem}

\begin{solution}
勾股定理：直角三角形两条直角边的平方和等于斜边的平方.

设 \(\triangle ABC\) 为直角三角形，且 \(\angle C=90^{\circ}\). 记
\(
BC=a,\quad CA=b,\quad AB=c.
\)
从点 \(C\) 向斜边 \(AB\) 作垂线，垂足为 \(H\).

因为 \(\angle AHC=\angle ACB=90^{\circ}\)，且 \(\angle CAH=\angle CAB\)，所以
\(
\triangle ACH\sim\triangle ABC.
\)
对应边的比例给出
\[
\frac{AC}{AB}=\frac{AH}{AC}
\]
从而
\[
AC^2=AB\cdot AH
\]

同理，\(\angle BHC=\angle BCA=90^{\circ}\)，且 \(\angle CBH=\angle CBA\)，所以
\(
\triangle BCH\sim\triangle BCA.
\)
因此
\[
\frac{BC}{BA}=\frac{BH}{BC}
\]
从而
\[
BC^2=AB\cdot BH
\]

将这两个等式相加，并利用 \(H\) 在线段 \(AB\) 上、故 \(AH+BH=AB\)，得到
\[
a^{2}+b^{2}=BC^{2}+CA^{2}=AB\cdot BH+AB\cdot AH=AB\left(AH+BH\right)=AB^{2}=c^{2}
\]

所以直角三角形两条直角边的平方和等于斜边的平方，勾股定理得证.
\end{solution}

\begin{problem}
外国船只除特许外不得进入离我海岸线 \(D\text{ 里}\) 以内的区域. 设 \(A\) 及 \(B\) 是我们的观测站，\(A\) 及 \(B\) 间的距离为 \(S\text{ 里}\)，海岸线是过 \(A\)，\(B\) 的直线，一外国船在 \(P\) 点，在 \(A\) 站测得 \(\angle BAP=\alpha\)，同时在 \(B\) 站测得 \(\angle ABP=\beta\). 则 \(\alpha\) 及 \(\beta\) 满足什么简单的三角函数值不等式，就应当向此未经特许的外国船发出警告，命令退出我海域？
\end{problem}

\begin{solution}
取海岸线 \(AB\) 为 \(x\) 轴，令
\(
A=(0,0),\quad B=(S,0).
\)
设外国船在海岸线同侧的坐标为
\(
P=(t,d),\qquad d>0.
\)
这里的 \(d\) 正是船到海岸线的距离，且不预先假定 \(P\) 的正射影落在线段 \(AB\) 内.

由
\(
\overrightarrow{AB}=(S,0),\quad \overrightarrow{AP}=(t,d)
\)
可得
\(\cos\alpha=\frac{t}{\sqrt{t^2+d^2}}\)，\(\sin\alpha=\frac{d}{\sqrt{t^2+d^2}}\)，\(\cot\alpha=\frac{t}{d}\).

在 \(B\) 点，观测基准方向是 \(\overrightarrow{BA}=(-S,0)\)，而
\(
\overrightarrow{BP}=(t-S,d).
\)
因此
\(\cos\beta=\frac{S-t}{\sqrt{(S-t)^2+d^2}}\)，\(\sin\beta=\frac{d}{\sqrt{(S-t)^2+d^2}}\)，\(\cot\beta=\frac{S-t}{d}\).
因此
\[
\cot\alpha+\cot\beta
=\frac{t}{d}+\frac{S-t}{d}
=\frac{S}{d}
\]
从而
\[
d=\frac{S}{\cot\alpha+\cot\beta}
\]

这里使用有向的水平分量，所以即使船的正射影落在线段 \(AB\) 外，公式仍然成立. 船进入距海岸线不超过 \(D\) 里的区域时应发出警告，即 \(d\leq D\). 又因为 \(d>0\)，所以 \(\cot\alpha+\cot\beta=S/d>0\)，代入上式得到
\[
\frac{S}{\cot\alpha+\cot\beta}\leq D
\quad\Longleftrightarrow\quad
\cot\alpha+\cot\beta\geq\frac{S}{D}
\]

因此应满足的三角函数值不等式为
\(
\cot\alpha+\cot\beta\geq\frac{S}{D}
\).
\end{solution}

\begin{problem}
设三棱锥 \(V - ABC\) 中，\(\angle AVB = \angle BVC = \angle CVA = 90^{\circ}\). 求证：\(\triangle ABC\) 是锐角三角形.
\end{problem}

\begin{solution}
设
\(
\bs a=\overrightarrow{VA},\quad \bs b=\overrightarrow{VB},\quad \bs c=\overrightarrow{VC}.
\)
因为三个棱间的夹角均为直角，所以
\[
\bs a\cdot\bs b=\bs b\cdot\bs c=\bs c\cdot\bs a=0
\]

在点 \(A\) 处，
\(
\overrightarrow{AB}=\bs b-\bs a,
\quad
\overrightarrow{AC}=\bs c-\bs a.
\)
于是
\[
\overrightarrow{AB}\cdot\overrightarrow{AC}
=(\bs b-\bs a)\cdot(\bs c-\bs a)
=\bs b\cdot\bs c-\bs b\cdot\bs a-\bs a\cdot\bs c+\lvert\bs a\rvert^2
=\lvert\bs a\rvert^2>0
\]
两条非零向量的内积为正，说明它们的夹角小于 \(90^{\circ}\)，所以 \(\angle BAC\) 为锐角.

在点 \(B\) 处，
\(
\overrightarrow{BA}=\bs a-\bs b,
\quad
\overrightarrow{BC}=\bs c-\bs b.
\)
因此
\[
\overrightarrow{BA}\cdot\overrightarrow{BC}
=(\bs a-\bs b)\cdot(\bs c-\bs b)
=\bs a\cdot\bs c-\bs a\cdot\bs b-\bs b\cdot\bs c+\lvert\bs b\rvert^2
=\lvert\bs b\rvert^2>0
\]
故 \(\angle ABC\) 为锐角.

在点 \(C\) 处，
\(
\overrightarrow{CA}=\bs a-\bs c,
\quad
\overrightarrow{CB}=\bs b-\bs c.
\)
所以
\[
\overrightarrow{CA}\cdot\overrightarrow{CB}
=(\bs a-\bs c)\cdot(\bs b-\bs c)
=\bs a\cdot\bs b-\bs a\cdot\bs c-\bs c\cdot\bs b+\lvert\bs c\rvert^2
=\lvert\bs c\rvert^2>0
\]
故 \(\angle ACB\) 也为锐角. 于是 \(\triangle ABC\) 的三个内角均为锐角，故 \(\triangle ABC\) 是锐角三角形.
\end{solution}

\begin{problem}
美国的物价从 1939 年的 \(100\) 增加到四十年后 1979 年的 \(500\)，如果每年物价增长率相同，问每年增长百分之几？（注意：\(x < 0.1\)，可用：\(\ln\left(1+x\right) \approx x\)，取 \(\lg 2 = 0.3\)，\(\ln 10 = 2.3\)）
\end{problem}

\begin{solution}
设每年的物价增长率为 \(x\). 因为物价逐年增长，所以 \(x>0\). 经过四十年，物价满足
\[
100(1+x)^{40}=500
\]
两边除以 \(100\)，得
\[
\left(1+x\right)^{40}=5
\]
两边取自然对数，得
\[
40\ln\left(1+x\right)=\ln5
\]
由换底关系 \(\ln a=(\lg a)\ln10\)，以及题给的 \(\lg2=0.3\)，可得
\[
\ln2=0.3\times2.3=0.69
\]
又因为 \(\lg5=\lg10-\lg2=1-0.3=0.7\)，所以
\[
\ln5=(\lg5)\ln10=0.7\times2.3=1.61
\]
因此
\[
\ln\left(1+x\right)=\frac{1.61}{40}=0.04025
\]
根据题设给出的近似式 \(\ln(1+x)\approx x\)，得到
\[
x\approx0.04025=4.025\%
\]
这个数小于题设给出的 \(0.1\)，与近似条件相符. 因此每年物价增长率约为 \(4.025\%\)，按题目数据通常取为约 \(4\%\).
\end{solution}

\begin{problem}
设 \(CEDF\) 是一个已知圆的内接矩形，过 \(D\) 作该圆的切线与 \(CE\) 的延长线相交于点 \(A\)，与 \(CF\) 的延长线相交于点 \(B\). 求证：\(\frac{BF}{AE} = \frac{BC^{3}}{AC^{3}}\).

\bitmapfigure[float=inline,width=5.8cm]{img/1979/national_science/q8_fig1.png}
\FloatBarrier
\end{problem}

\begin{solution}
因为 \(CEDF\) 是矩形，所以 \(CE\perp CF\). 又因为 \(A\)，\(C\)，\(E\) 共线、\(B\)，\(C\)，\(F\) 共线，所以
\[
\angle ACB=90^{\circ}
\]
矩形的两个对角顶点为 \(C\)，\(D\)，因此对角线 \(CD\) 是该圆的直径. 圆在 \(D\) 点的切线 \(AB\) 垂直于半径 \(OD\)，而 \(O\)，\(C\)，\(D\) 共线，故
\[
CD\perp AB
\]
于是 \(CD\) 是直角三角形 \(ACB\) 斜边 \(AB\) 上的高，\(\triangle ACD\) 和 \(\triangle BCD\) 都是直角三角形.

由 \(\triangle ACD\sim\triangle ABC\)，有
\[
\frac{AD}{AC}=\frac{AC}{AB}
\]
所以
\[
AD=\frac{AC^2}{AB}
\]
同理，由 \(\triangle BCD\sim\triangle BCA\)，有
\[
\frac{BD}{BC}=\frac{BC}{BA}
\]
所以
\[
BD=\frac{BC^2}{AB}
\]
两式相除，得到
\[
\frac{BD}{AD}=\frac{BC^2}{AC^2}
\]

再由切线与割线的点幂定理，在点 \(A\) 处切线段为 \(AD\)，割线依次交圆于 \(E\)，\(C\)，故
\[
AD^2=AE\cdot AC
\]
在点 \(B\) 处切线段为 \(BD\)，割线依次交圆于 \(F\)，\(C\)，故
\[
BD^2=BF\cdot BC
\]
由此
\(AE=\frac{AD^2}{AC}\)，\(BF=\frac{BD^2}{BC}\).
于是
\[
\frac{BF}{AE}
=\frac{BD^2/BC}{AD^2/AC}
=\left(\frac{BD}{AD}\right)^2\frac{AC}{BC}
=\left(\frac{BC^2}{AC^2}\right)^2\frac{AC}{BC}
=\frac{BC^3}{AC^3}
\]
故命题得证.
\end{solution}

\begin{problem}
试问数列 \(\lg 100\)，\(\lg \left(100\sin \frac{\pi}{4}\right)\)，\(\lg \left(100\sin^{2}\frac{\pi}{4}\right)\)，\(\cdots\)，\(\lg \left(100\sin^{n-1}\frac{\pi}{4}\right)\) 前多少项的和的值最大？并求这最大值. （\(\lg 2 = 0.301\)）
\end{problem}

\begin{solution}
记第 \(k\) 项为 \(a_k\)，其中 \(k\geq1\). 由于
\[
\sin\frac{\pi}{4}=\frac{\sqrt{2}}{2}=2^{-1/2}
\]
且 \(\lg100=2\)，所以
\[
a_k=\lg\left(100\left(\sin\frac{\pi}{4}\right)^{k-1}\right)
=\lg\left(100\cdot2^{-(k-1)/2}\right)
=2-\frac{k-1}{2}\lg2
=2-0.1505(k-1)
\]

所以 \(a_k\) 随 \(k\) 增大而严格递减. 记前 \(k\) 项的和为 \(S_k\)，则
\[
S_{k+1}-S_k=a_{k+1}=2-0.1505k
\]
考察最大值附近的两项：
\[
a_{14}=2-0.1505\times13=0.0435>0
\]
\[
a_{15}=2-0.1505\times14=-0.107<0
\]
因此 \(S_{14}-S_{13}=a_{14}>0\)，而 \(S_{15}-S_{14}=a_{15}<0\). 又因为所有 \(a_k\) 递减，所以前 \(14\) 项的加法过程一直增加到 \(S_{14}\)，从第 \(15\) 项起继续加项都会使和减小，故最大值在前 \(14\) 项处取得.

前 \(14\) 项构成首项为 \(2\)、第 \(14\) 项为 \(0.0435\) 的等差数列，故其和为
\[
S_{14}=\frac{14(2+0.0435)}{2}=14.3045
\]
所以前 \(14\) 项的和最大，最大值为 \(14.3045\)（按给定近似数据约为 \(14.30\)）.
\end{solution}

\begin{problem}
设等腰 \(\triangle OAB\) 的顶角为 \(2\theta\)，高为 \(h\).

\begin{enumerate}
\item \(\triangle OAB\) 内有一动点 \(P\)，到三边 \(OA\)，\(OB\)，\(AB\) 的距离分别为 \(\lvert PD\rvert\)，\(\lvert PF\rvert\)，\(\lvert PE\rvert\)，并且满足关系 \(\lvert PD\rvert \cdot \lvert PF\rvert = \lvert PE\rvert^2\). 求 \(P\) 点的轨迹；
\item 在上述轨迹中定出点 \(P\) 的坐标，使得 \(\lvert PD\rvert + \lvert PE\rvert = \lvert PF\rvert\).
\end{enumerate}
\end{problem}

\begin{solution}
\begin{enumerate}
\item 取 \(O\) 为原点，以高从 \(O\) 指向边 \(AB\) 的方向为 \(x\) 轴正向，以垂直于高的方向为 \(y\) 轴正向，并取 \(A\) 位于 \(x\) 轴上方. 由于顶角为 \(2\theta\)，有 \(0<\theta<90^{\circ}\)，于是
\(A=(h,h\tan\theta)\)，\(B=(h,-h\tan\theta)\).

设 \(P=(x,y)\). 因为 \(P\) 在三角形内部，所以
\(0<x<h\)，\(-x\tan\theta<y<x\tan\theta\).
三边所在直线分别为
\(OA: y=x\tan\theta\)，\(OB: y=-x\tan\theta\)，\(AB: x=h\).
由点到直线的距离公式，且利用 \(\sqrt{1+\tan^2\theta}=1/\cos\theta\)，可得
\[
\lvert PD\rvert
=\frac{x\tan\theta-y}{\sqrt{1+\tan^2\theta}}
=x\sin\theta-y\cos\theta
\]
\[
\lvert PF\rvert
=\frac{x\tan\theta+y}{\sqrt{1+\tan^2\theta}}
=x\sin\theta+y\cos\theta
\]
并且
\[
\lvert PE\rvert=h-x
\]
这些量在三角形内部均为正.

代入 \(\lvert PD\rvert\lvert PF\rvert=\lvert PE\rvert^2\)，得到
\[
x^2\sin^2\theta-y^2\cos^2\theta=(h-x)^2
\]
展开右边并移项，得
\(c^2\left(x^2+y^2\right)-2hx+h^2=0\)，\(c=\cos\theta\).
将 \(c\) 换回 \(\cos\theta\)，配方得
\[
\left(x-\frac{h}{\cos^2\theta}\right)^2+y^2
=\left(\frac{h\sin\theta}{\cos^2\theta}\right)^2
\]
因此，满足距离乘积条件的点在上述圆上；但题目还要求 \(P\) 位于三角形内部，所以实际轨迹是该圆满足
\(0<x<h\)，\(-x\tan\theta<y<x\tan\theta\)
的那一段（边界因 \(P\) 在三角形内部而不取）.

\item 由 \(\lvert PD\rvert+\lvert PE\rvert=\lvert PF\rvert\)，有
\[
x\sin\theta-y\cos\theta+h-x=x\sin\theta+y\cos\theta
\]
从而
\[
h-x=2y\cos\theta
\]
因为 \(0<x<h\) 且 \(\cos\theta>0\)，所以 \(y>0\).

把 \(y\cos\theta=(h-x)/2\) 代入乘积条件的等价形式，得到
\[
x^2\sin^2\theta-\frac{(h-x)^2}{4}=(h-x)^2
\]
整理为
\[
4x^2\sin^2\theta=5(h-x)^2
\]
由于 \(x>0\)、\(\sin\theta>0\) 且 \(h-x>0\)，等式两边取正平方根，得到
\[
2x\sin\theta=\sqrt{5}\left(h-x\right)
\]
解得
\[
x=\frac{\sqrt{5}h}{\sqrt{5}+2\sin\theta}
\]
再由 \(h-x=2y\cos\theta\)，得到
\[
y=\frac{h-x}{2\cos\theta}=\frac{h\tan\theta}{\sqrt{5}+2\sin\theta}
\]
所以所求点的坐标为
\[
P=\left(\frac{\sqrt{5}h}{\sqrt{5}+2\sin\theta},
\frac{h\tan\theta}{\sqrt{5}+2\sin\theta}\right)
\]

最后验证位置条件：分母大于 \(\sqrt5\)，所以 \(0<x<h\)，并且
\[
\frac{y}{x\tan\theta}=\frac{1}{\sqrt5}<1
\]
由于 \(y>0\)，故 \(-x\tan\theta<y<x\tan\theta\)，该点确实在三角形内部；代回上面的等式可知它也在第一问所得轨迹上.
\end{enumerate}
\end{solution}
